\documentclass[11pt]{article}
\PassOptionsToPackage{gray}{xcolor}
\usepackage[a4paper,margin=22mm]{geometry}
\usepackage{graphicx}
\usepackage{booktabs}
\usepackage{xcolor}
\usepackage{fontspec}
\usepackage[unicode,colorlinks=true,linkcolor=teal,urlcolor=teal]{hyperref}

\setmainfont[BoldFont=lmroman10-bold.otf]{lmroman10-regular.otf}
\definecolor{ReportInk}{HTML}{172B33}
\definecolor{ReportAccent}{HTML}{286C68}
\definecolor{ReportTint}{HTML}{EAF3F2}

\setlength{\parindent}{0pt}
\setlength{\parskip}{5pt}
\linespread{1.08}
\renewcommand{\arraystretch}{1.26}
\setlength{\tabcolsep}{5pt}

\newcommand{\reportlabel}[1]{{\large\color{ReportAccent}\bfseries #1}}
\newcommand{\reporttitle}[1]{%
  {\color{ReportInk}\Huge\bfseries #1\par}%
  \vspace{5mm}{\color{ReportAccent}\hrule height 1.3pt}%
}
\newcommand{\reportfigure}[3]{%
  \begin{center}
    {\bfseries #1}\par\smallskip
    \includegraphics[width=0.90\linewidth]{#2}\par\smallskip
    \begin{minipage}{0.9\linewidth}\small #3\end{minipage}
  \end{center}%
}

% Generated by python3 -m src.phase1_tables
% Localized reports can override these labels and number formats before loading this file.
\providecommand{\ReportNumber}[1]{#1}
\providecommand{\ReportRange}[2]{#1--#2}
\providecommand{\ReportMissing}{Not available}
\providecommand{\TableUDS}{UDS}
\providecommand{\TableMRI}{MRI}
\providecommand{\TableParticipant}{Participant}
\providecommand{\TableMatrix}{Adjacency matrix}
\providecommand{\TableRegion}{Brain region}
\providecommand{\TableAnimals}{Animal-naming score}
\providecommand{\TableAnimalsShort}{Animal-naming score}
\providecommand{\TableVegetables}{Vegetable-naming score}
\providecommand{\TableAge}{Age at visit}
\providecommand{\TableSex}{Recorded sex}
\providecommand{\TableEducation}{Years of education}
\providecommand{\TableThickness}{Left cortical thickness}
\providecommand{\TableMatched}{Matched participant table}
\providecommand{\TableStatusAndMRI}{Usable cognitive status and thickness}
\providecommand{\TablePairedAll}{Paired score and thickness, all groups}
\providecommand{\TablePairedNormal}{Paired score and thickness, normal cognition}
\providecommand{\TableModelAll}{Supplementary model, all groups}
\providecommand{\TableModelNormal}{Primary model, normal cognition}
\providecommand{\TableThicknessMM}{Left thickness (mm)}
\providecommand{\TableAgeYears}{Age (years)}
\providecommand{\TableEducationYears}{Education (years)}
\providecommand{\TableNormal}{Normal cognition}
\providecommand{\TableImpairedWithoutMCI}{Impaired without MCI}
\providecommand{\TableMCI}{MCI}
\providecommand{\TableDementia}{Dementia}
\newcommand{\NMatched}{\ReportNumber{1,237}}
\newcommand{\NNormal}{\ReportNumber{847}}
\newcommand{\NPairedNormal}{\ReportNumber{843}}
\newcommand{\NModelNormal}{\ReportNumber{842}}
\newcommand{\NOnlyUDS}{\ReportNumber{0}}
\newcommand{\NOnlyMRI}{\ReportNumber{0}}
\newcommand{\NStatusEqual}{\ReportNumber{1,237}}
\newcommand{\NNoScoreNormal}{\ReportNumber{4}}
\newcommand{\NNoEducationNormal}{\ReportNumber{1}}
\newcommand{\NormalScoreMedian}{\ReportNumber{22}}
\newcommand{\NormalThicknessMedian}{\ReportNumber{2.338}}
\newcommand{\NormalScoreMin}{\ReportNumber{7}}
\newcommand{\NormalScoreMax}{\ReportNumber{43}}
\newcommand{\NormalThicknessMin}{\ReportNumber{1.474}}
\newcommand{\NormalThicknessMax}{\ReportNumber{2.666}}
\newcommand{\NormalScoreFlags}{\ReportNumber{4}}
\newcommand{\NormalThicknessFlags}{\ReportNumber{15}}
\newcommand{\NormalAgeMean}{\ReportNumber{67.27}}
\newcommand{\MCIAgeMean}{\ReportNumber{70.83}}
\newcommand{\NWomen}{\ReportNumber{740}}
\newcommand{\NSex}{\ReportNumber{1,237}}
\newcommand{\SourceRows}{%
\TableUDS{} & \ReportNumber{1,237} & \ReportNumber{180} & \TableParticipant{} \\
\TableMRI{} & \ReportNumber{1,237} & \ReportNumber{126} & \TableParticipant{} \\
\TableMatrix{} & \ReportNumber{62} & \ReportNumber{63} & \TableRegion{} \\
}
\newcommand{\QualityRows}{%
\TableAnimals{} & \ReportNumber{1,225} & \ReportNumber{12} & \ReportNumber{0.97}\% \\
\TableVegetables{} & \ReportNumber{1,222} & \ReportNumber{15} & \ReportNumber{1.21}\% \\
\TableAge{} & \ReportNumber{1,237} & \ReportNumber{0} & \ReportNumber{0.00}\% \\
\TableSex{} & \ReportNumber{1,237} & \ReportNumber{0} & \ReportNumber{0.00}\% \\
\TableEducation{} & \ReportNumber{1,235} & \ReportNumber{2} & \ReportNumber{0.16}\% \\
\TableThickness{} & \ReportNumber{1,237} & \ReportNumber{0} & \ReportNumber{0.00}\% \\
}
\newcommand{\FlowRows}{%
\TableMatched{} & \ReportNumber{1,237} \\
\TableStatusAndMRI{} & \ReportNumber{1,237} \\
\TablePairedAll{} & \ReportNumber{1,225} \\
\TablePairedNormal{} & \ReportNumber{843} \\
\TableModelAll{} & \ReportNumber{1,223} \\
\TableModelNormal{} & \ReportNumber{842} \\
}
\newcommand{\DescriptiveRows}{%
\TableAnimalsShort{} & \ReportNumber{843} & \ReportNumber{22.32} & \ReportNumber{5.43} & \ReportNumber{22} & \ReportRange{\ReportNumber{18}}{\ReportNumber{26}} & \ReportNumber{8} \\
\TableThicknessMM{} & \ReportNumber{847} & \ReportNumber{2.336} & \ReportNumber{0.121} & \ReportNumber{2.338} & \ReportRange{\ReportNumber{2.269}}{\ReportNumber{2.417}} & \ReportNumber{0.148} \\
\TableAgeYears{} & \ReportNumber{847} & \ReportNumber{67.27} & \ReportNumber{7.92} & \ReportNumber{67} & \ReportRange{\ReportNumber{63}}{\ReportNumber{72}} & \ReportNumber{9} \\
\TableEducationYears{} & \ReportNumber{846} & \ReportNumber{16.38} & \ReportNumber{2.61} & \ReportNumber{16} & \ReportRange{\ReportNumber{15}}{\ReportNumber{18}} & \ReportNumber{3} \\
}
\newcommand{\GroupRows}{%
\TableNormal{} & \ReportNumber{847} & \ReportNumber{843} & \ReportNumber{22.32} & \ReportNumber{847} & \ReportNumber{2.336} \\
\TableImpairedWithoutMCI{} & \ReportNumber{7} & \ReportNumber{7} & \ReportNumber{22.00} & \ReportNumber{7} & \ReportNumber{2.314} \\
\TableMCI{} & \ReportNumber{252} & \ReportNumber{249} & \ReportNumber{17.83} & \ReportNumber{252} & \ReportNumber{2.258} \\
\TableDementia{} & \ReportNumber{131} & \ReportNumber{126} & \ReportNumber{12.27} & \ReportNumber{131} & \ReportNumber{2.154} \\
}

\graphicspath{{figures/phase1/}}
\hypersetup{pdftitle={Semantic verbal fluency and left inferior parietal cortical thickness}}

\hypersetup{allcolors=black,linkbordercolor={0 0 0},urlbordercolor={0 0 0}}

\begin{document}
\reportlabel{Statistics course project, Track C (NACC)}
\vspace{9mm}

\reporttitle{Semantic verbal fluency and\\left inferior parietal cortical thickness}
\vspace{3mm}
{\Large Phase 1 | Study design, data, and descriptive statistics\par}
\vspace{9mm}

{\large\bfseries Research question}\par
Among participants with normal cognition, how many more or fewer animals
are named correctly in 60 seconds per 0.1 mm greater left inferior
parietal cortical thickness, after accounting for age, sex, and years of
education?

\vspace{5mm}
\noindent\colorbox{ReportTint}{%
  \parbox{\dimexpr\linewidth-2\fboxsep\relax}{%
    \vspace{2mm}
    The teaching subset contains clinical and MRI records for \NMatched{}
    matched participants. Of the \NNormal{} classified as cognitively normal,
    \NPairedNormal{} have usable scores and thickness measurements, and
    \NModelNormal{} have complete data for the adjusted model.
    \vspace{2mm}
  }%
}

\vspace{5mm}
{\large\bfseries Phase 1 summary}\par
We linked the clinical and MRI files by participant ID and documented how
we handled special codes, out-of-range values, and fractional scores.
This report describes the distribution of cognitive status, summarizes
the main variables, and compares the recorded groups. Among participants
with normal cognition, the median score was \NormalScoreMedian{} correct
responses and the median left cortical thickness was
\NormalThicknessMedian{} mm. Later phases will estimate the within-group
association and account for the covariates.

\vfill
{\small Autumn 1405 (Solar Hijri)\par}
\clearpage

\section{Data source, structure, and variables}
This secondary observational analysis uses a teaching subset from the
National Alzheimer's Coordinating Center (NACC), supplied by the course
team. The Uniform Data Set (UDS) contains clinical assessments and
cognitive tests. The MRI file contains measurements derived from
structural imaging. The file columns do not identify the date of
receipt, release version, or collection period for this extract.

\begin{center}\small
\begin{tabular}{@{}lrrl@{}}
\toprule
Supplied file & Rows & Columns & Unit of observation \\
\midrule
\SourceRows
\bottomrule
\end{tabular}
\end{center}

The matrix column count includes one column for row labels. The 62
region names match the MRI columns. The MRI file contains one cortical
thickness measure and one gray matter volume measure for each region.
The matrix is documented in the data inventory. The primary model uses
the left-side thickness measurement.

\subsection*{Main variables}
\begin{center}\small
\begin{tabular}{@{}lp{0.36\linewidth}l@{}}
\toprule
Variable & Definition and scale & Role \\
\midrule
ANIMALS & Correct responses in 60 seconds, discrete & Outcome \\
LH\_INFERIORPARIETAL\_AVGTH & Left cortical thickness, continuous, in mm & Explanatory variable \\
NACCAGE & Age at the visit, quantitative, in years & Covariate \\
SEX & Recorded sex, nominal & Covariate \\
EDUC & Years of education, discrete & Covariate \\
NACCUDSD & Recorded cognitive status, categorical & Defines the groups \\
\bottomrule
\end{tabular}
\end{center}

The UDS inventory groups columns by content, including demographics,
medical history and medications, everyday functioning, cognitive tests,
and diagnosis. The audit output lists every column and the notebook
shows the number of columns in each category.
\clearpage

\section{Data audit and sample flow}
Each participant ID (NACCID) appeared once in each file. We joined the
records one-to-one on that identifier. There were \NOnlyUDS{} UDS-only
and \NOnlyMRI{} MRI-only IDs. Cognitive status agreed across files for
\NStatusEqual{} matched participants. Group labels come from UDS.

\begin{center}\small
\begin{tabular}{@{}lrrr@{}}
\toprule
Variable & Usable & Unusable & Unusable (\%) \\
\midrule
\QualityRows
\bottomrule
\end{tabular}
\end{center}

ANIMALS and VEG are whole-number counts. For ANIMALS, we kept the eleven
noncompletion codes and the value of 21.6 in the raw data and excluded
them from the usable scores. We left 21.6 as supplied because its origin
is unknown. Education code 99 indicates an unknown value. All recorded
left cortical thickness measurements were finite and positive. The
subset provides no information on scan-level quality control.

\begin{center}\small
\begin{tabular}{@{}lr@{}}
\toprule
Analysis sample & Participants \\
\midrule
\FlowRows
\bottomrule
\end{tabular}
\end{center}

Among participants with normal cognition, \NNoScoreNormal{} lacked a
usable animal-naming score and \NNoEducationNormal{} additional
participant had unknown education. Every UDS participant in the
supplied files has an MRI record, so there is no group without MRI for
the descriptive comparison.
\clearpage

\section{Descriptive statistics}
The table uses all usable values for each variable in the
normal-cognition group. The count in each row is the denominator for
that variable. We calculated quartiles by linear interpolation and
sample standard deviations with denominator n-1.

The mean and standard deviation describe center and spread in the
original units, making comparisons across groups easier. The median
and interquartile range (IQR) describe the center and middle half of
the distribution with less sensitivity to extreme values. The mean
and median of the animal-naming score are close. Thickness has a tail
toward lower values, so the median and IQR are useful summaries of
its central portion.

\begin{center}\small
\begin{tabular}{@{}lrrrrrr@{}}
\toprule
Variable & Usable & Mean & SD & Median & Q1--Q3 & IQR \\
\midrule
\DescriptiveRows
\bottomrule
\end{tabular}
\end{center}

In the normal-cognition group, scores ranged from \NormalScoreMin{} to
\NormalScoreMax{} correct responses and thickness ranged from
\NormalThicknessMin{} to \NormalThicknessMax{} mm. The IQR rule flagged
\NormalScoreFlags{} scores and \NormalThicknessFlags{} thickness
measurements for review. All remained in the summaries.

\subsection*{Comparison across cognitive status groups}
The table reports a separate denominator for each measure. Means are
given in correct responses and millimeters, respectively.

\begin{center}\small
\begin{tabular}{@{}lrrrrr@{}}
\toprule
Cognitive status & Group total & Usable scores & Mean score &
\shortstack[r]{Usable\\thickness} & \shortstack[r]{Mean\\thickness} \\
\midrule
\GroupRows
\bottomrule
\end{tabular}
\end{center}

The normal-cognition group averaged about 22 correct responses. Mean
scores and thickness were lower in the MCI and dementia groups. The
group classified as impaired without MCI contains seven people, so a
few observations can noticeably affect its summaries. These
comparisons are unadjusted. Mean age was \NormalAgeMean{} years in the
normal-cognition group and \MCIAgeMean{} years in the MCI group.
Interpretation is limited to the composition of these groups.

Across the full sample, \NWomen{} of the \NSex{} participants with a
usable sex code were recorded as female. Figure 3 shows the counts and
percentages for all cognitive status groups. The notebook also gives
observed and cumulative frequencies for scores and years of
education, measures of spread, and the analysis sample denominators.
\clearpage

\section{Distributions in the normal-cognition group}
\reportfigure{Figure 1. Animal-naming scores}{01_animal_naming.pdf}{%
  Each bar counts participants with a given whole-number score. The
  middle half of observed scores fell between 18 and 26 correct responses.}
\reportfigure{Figure 2. Left inferior parietal cortical thickness}{02_left_parietal_thickness.pdf}{%
  The bins are 0.05 mm wide. Measurements cluster near 2.3 mm,
  with a tail toward lower values.}
\clearpage

\section{Cognitive status frequencies and group comparisons}
\reportfigure{Figure 3. Cognitive status in the study sample}{03_cognitive_status.pdf}{%
  Each percentage uses all matched participants as its denominator.
  These percentages describe the composition of the teaching subset.}
\reportfigure{Figure 4. Left cortical thickness by cognitive status}{04_thickness_by_status.pdf}{%
  The line inside each box marks the median, and the box edges mark the
  first and third quartiles. Circles mark values outside the 1.5-IQR
  fences. These values are included in the descriptive summaries.}
\clearpage

\section{Conclusion and limitations}
\subsection*{Scope of the final question}
The primary question concerns the association between animal-naming
performance and left inferior parietal cortical thickness among
participants with normal cognition. Age, sex, and years of education
are the model covariates. Both measures are available at the
participant level and can be linked one-to-one, leaving \NModelNormal{}
participants with complete data for the primary model. Comparisons
across the full set of cognitive status groups describe the sample
and provide context for the later analysis.

\subsection*{Data limitations}
\begin{itemize}
\item Participants entered the teaching subset through random selection. The sampling frame and eligibility criteria are undocumented.
Interpretation is limited to the supplied subset and the defined
analysis sample.
\item Each participant has one record per source in the teaching
subset. The files do not specify the number of visits in the source
database, the scan date, or the interval between the scan and the test.
\item The preparation version, the origin of fractional scores, and
scan-level quality control details are unknown. VEG and a bilateral
thickness definition were specified in advance as sensitivity analysis
options.
\end{itemize}

\subsection*{References}
\begin{enumerate}
\item Official statistics course project guide, Track C, pp. 4--8 and 18--19.
\item \href{https://www.naccdata.org/the-nacc-researchers-guide/}{NACC Researcher's Guide}, data structure and linkage of imaging to clinical visits.
\item \href{https://r2d2.kumc.edu/ADC/Protocols/UDS2_protocol/ResearchDataDict_2015Mar.pdf\#page=205}{UDS data dictionary}, score definitions and special codes.
\item \href{https://files.alz.washington.edu/documentation/SCAN-MRI-Imaging-RDD.pdf\#page=17}{MRI data dictionary}, inferior parietal cortical thickness and its unit.
\item \href{https://pmc.ncbi.nlm.nih.gov/articles/PMC4207059/}{Eastman et al. (2013)}, cortical thickness and semantic verbal fluency.
\end{enumerate}
\end{document}
